\documentclass[11pt]{article}
\usepackage[T1]{fontenc}
\usepackage[utf8]{inputenc}
\usepackage{lmodern}
\usepackage{amsmath,amssymb}
\usepackage{geometry}
\usepackage{hyperref}
\geometry{margin=1in}

\title{Schwinger-Model Pion Form Factor in \texttt{jaxQFT}}
\author{}
\date{\today}

\begin{document}
\maketitle

\section{Scope}
This note documents the DD-independent form-factor workflow now implemented for
the two-dimensional U(1) Wilson-fermion Schwinger-model example in
\texttt{jaxQFT}. The target observable is the pion matrix element of the
isovector vector current, measured with dense all-to-all propagators and
analyzed with autocorrelation-aware blocked jackknife fits.

\section{Current and Interpolator}
The source and sink interpolators are the local pseudoscalar operators
\begin{equation}
  P^+(x) = \bar d(x) \gamma_5 u(x),
  \qquad
  P^-(x) = \bar u(x) \gamma_5 d(x),
\end{equation}
and the current insertion is either the local isovector vector current
\begin{equation}
  V^{(3)}_\mu(x) = \bar\psi(x)\gamma_\mu \frac{\tau^3}{2}\psi(x).
\end{equation}
or the exactly conserved Wilson point-split current
\begin{equation}
  J^{(3)}_\mu\!\left(x+\frac{\hat\mu}{2}\right)
  =
  \frac{1}{2}\Big[
    \bar\psi(x+\hat\mu)\,U^\dagger_\mu(x)\,(r+\gamma_\mu)\,\frac{\tau^3}{2}\,\psi(x)
    -
    \bar\psi(x)\,U_\mu(x)\,(r-\gamma_\mu)\,\frac{\tau^3}{2}\,\psi(x+\hat\mu)
  \Big].
\end{equation}
For the charged-pion connected correlator, the $u$-line and $d$-line current
insertions are equal, so the connected three-point function is represented by a
single quark-line triangle contraction.

\section{Dense All-to-All Construction}
Let $S(x,y)$ denote the full quark propagator. For source point $y$, current
insertion point $z$, and sink point $x$, the connected pion three-point kernel
implemented in the code is
\begin{equation}
  K_\mu(x,z,y)
  =
  \mathrm{Tr}\left[
    \gamma_5 S(x,z) \gamma_\mu S(z,y) \gamma_5 S(y,x)
  \right].
\end{equation}
The inline measurement uses the dense inverse of the Wilson operator, reshaped
into timeslice blocks, so that for each chosen source time $t_{\rm src}$, sink
separation $t_{\rm sep}$, and insertion time $\tau$, the spatial sums are done
exactly. For the conserved current, the insertion is evaluated as the sum of
its forward and backward point-split pieces with the same dense all-to-all data.

\section{Momentum Projection}
With source time $t_{\rm src}$, sink time $t_{\rm src}+t_{\rm sep}$, and
insertion time $t_{\rm src}+\tau$, the measured correlator is
\begin{equation}
  C_{3,\mu}(\tau,t_{\rm sep};p_f,p_i)
  =
  \frac{1}{N_{\rm src}}
  \sum_{t_{\rm src},y,x,z}
  e^{-ip_f x}
  e^{+i(p_f-p_i)z}
  e^{+ip_i y}
  K_\mu(x,z,y),
\end{equation}
where $N_{\rm src}$ is the number of source sites used in the translation
average. In the current implementation:
\begin{itemize}
  \item all spatial source positions on each selected source timeslice are used,
  \item the source-time average is configurable,
  \item the measurement currently targets 2D lattices $(L_x,L_t)$ with momentum
  along the single spatial direction.
\end{itemize}

\section{Practical Run Setup}
The production card added for the 16\,$\times$\,32 quenched example is
\begin{verbatim}
runs/U1-example/mcmc_u1_quenched_beta8_k0p2530_16x32_pion_form_factor.toml
\end{verbatim}
with the following choices:
\begin{itemize}
  \item quenched gauge evolution at $\beta=8.0$,
  \item dense, translation-averaged pion two-point functions for
  $p\in\{-2,-1,0,1,2\}$,
  \item dense pion three-point functions in Breit kinematics
  $(p_f,p_i)=(p,-p)$ for $p=0,1,2$,
  \item temporal current only in the example card ($\mu=1$ in this 2D convention),
  \item all source times, and all source-sink separations
  $t_{\rm sep}=4,5,\ldots,16$.
\end{itemize}
The example card currently defaults to \texttt{current\_kind = "local"}.
Setting \texttt{current\_kind = "conserved"} switches to the point-split
current while keeping the rest of the workflow unchanged.
This is a deliberate compromise: it exploits full translation invariance while
keeping the number of stored three-point channels small enough for a practical
checkpoint size.

\section{Output Layout}
The three-point measurement is stored under the inline measurement name
\texttt{pion\_3pt\_vector} with flat scalar keys of the form
\begin{verbatim}
c3_mu{mu}_pf{pf}_pi{pi}_tsep{tsep}_tau{tau}_re
c3_mu{mu}_pf{pf}_pi{pi}_tsep{tsep}_tau{tau}_im
\end{verbatim}
This matches the single-channel extractor already present in
\texttt{scripts/mcmc/analyze\_pion\_vector\_3pt.py}.

\section{Analysis Path}
Two analysis scripts are now available.

\subsection{Single channel}
\texttt{scripts/mcmc/analyze\_pion\_vector\_3pt.py} extracts one selected
channel $(\mu,p_i,p_f,t_{\rm sep})$ at a time. It fits the two-point
functions, removes the overlap factors, performs a correlated plateau fit, and
divides by the kinematic prefactor $(p_f+p_i)_\mu$ whenever that prefactor is
nonzero. In the present 2D convention this means
\begin{itemize}
  \item $\mu=1$ (temporal): divide by $E_f+E_i$,
  \item $\mu=0$ (spatial): divide by $p_f+p_i$.
\end{itemize}
The analysis also has to choose the correct complex component of the reduced
matrix element. In the present Euclidean convention this is
\begin{itemize}
  \item $\mathrm{Re}\,M_{\rm eff}$ for the temporal current,
  \item $-\mathrm{Im}\,M_{\rm eff}$ for the spatial current.
\end{itemize}
This is now the default \texttt{--component auto} behavior in both analysis
scripts, and it is the convention used when combining $\mu=0$ and $\mu=1$
channels into a common fixed-$Q^2$ average.
Therefore the spatial current is useful for non-Breit kinematics, while it
vanishes identically in the Breit frame.

The script reports two estimators for the same matrix element:
\begin{enumerate}
  \item a ratio estimator, where the source/sink overlap factors cancel
  against appropriate two-point functions,
  \item a direct estimator, where the overlap factors are reconstructed from
  separate two-point fits through
  $Z_\pi(p)=\sqrt{2E_\pi(p)A_\pi(p)}$.
\end{enumerate}
In both cases the fitted pion energies are also used to determine the momentum
transfer
\begin{equation}
  Q^2 = (E_f-E_i)^2 + (p_f-p_i)^2
\end{equation}
in Euclidean conventions. This is not approximated by the spatial momentum
transfer alone.

\subsection{Batch form-factor scan}
\texttt{scripts/mcmc/analyze\_pion\_form\_factor.py} scans all matching
channels in a checkpoint, runs the same blocked-jackknife extraction for each
one, and summarizes the results in a JSON table together with a summary PNG.
The reduction is done in two stages:
\begin{enumerate}
  \item First, all channels that share the same $(Q^2,t_{\rm sep})$ are
  combined with a covariance-weighted average. This is where different current
  components $\mu$ and different momentum routings contribute to the same
  matrix element after the appropriate kinematic prefactors have been removed.
  \item Second, the resulting fixed-$Q^2$ data are analyzed as a function of
  $t_{\rm sep}$, and an optional correlated constant fit can be used as a
  \emph{late-$t_{\rm sep}$ diagnostic}. This is useful for checking whether the
  large-$t_{\rm sep}$ points are mutually consistent, but it should not be
  confused with a full excited-state analysis.
\end{enumerate}
For the default 16\,$\times$\,32 card this means a scan over the available
Breit-frame momentum transfers and all stored source-sink separations. If the
measurement was done with \texttt{pair\_mode = "all"} and both current
components were stored, the batch script can additionally be run with
\texttt{--all-mus}. In that mode it computes correlated averages of all channels
that share the same $(Q^2,t_{\rm sep})$, so the spatial current contributes
extra statistics whenever $p_f+p_i\neq 0$.
For charge-normalization studies, the batch script also supports
\texttt{--pair-mode q2zero} (aliases \texttt{equal}, \texttt{diagonal}),
which restricts the scan to the equal-momentum channels $p_f=p_i$.

\section{Statistics}
The analysis is blocked and IAT-aware:
\begin{enumerate}
  \item two-point and three-point samples are aligned on common MCMC steps,
  \item a block size is chosen automatically from integrated autocorrelation
  times unless overridden,
  \item pion energies and overlap factors are fit on blocked data,
  \item the $\tau$ plateau in each 3pt channel is chosen with correlated
  constant fits,
  \item channels with the same $(Q^2,t_{\rm sep})$ are combined with their full
  jackknife covariance,
  \item optionally, the remaining $t_{\rm sep}$ dependence at fixed $Q^2$ is
  summarized with a second correlated constant fit as a late-$t_{\rm sep}$
  consistency check,
  \item matrix elements and form factors are quoted from blocked jackknife
  replicas.
\end{enumerate}

\section{Current Limitation}
Both current choices are now available through
\texttt{current\_kind = "local"|"conserved"} in the inline measurement block.
For the conserved current one should use $Z_V=1$ in the analysis scripts. For
the local current the same analysis path applies, but a separate multiplicative
renormalization factor is still needed. The conserved-current sign convention
now matches the local-current path, and the analysis scripts recover
$F_\pi(Q^2=0)=1$ to machine precision on a synthetic single-state checkpoint.
The remaining caveat is practical rather than formal: a tiny hot-start dataset
such as $4\times 8$ with $t_{\rm sep}=4$ is not a meaningful charge-normalization
test because wrap-around and excited-state contamination are still large there.
A proper physics validation should therefore be done on a thermalized ensemble
with several source-sink separations.

Another limitation should be stated explicitly: the current code does
\emph{not} yet perform a single global simultaneous fit of the full two-point
and three-point data set. The implemented workflow is instead
\begin{enumerate}
  \item separate blocked fits of the required two-point functions,
  \item ratio/direct-overlap reconstruction of each 3pt matrix element,
  \item correlated reduction over channels and $t_{\rm sep}$.
\end{enumerate}
This is already a statistically controlled extraction strategy, but it is not
the same as a joint $(C_2,C_3)$ simultaneous fit with common energies,
overlaps, and matrix elements.

Relatedly, the current batch reduction does not yet model the residual excited
state contamination in the $t_{\rm sep}$ dependence of the reduced form factor.
At fixed $Q^2$ one should expect
\begin{equation}
  F_{\rm eff}(Q^2,t_{\rm sep})
  =
  F_\pi(Q^2)
  +
  \mathcal{O}(e^{-\Delta E\, t_{\rm sep}})
\end{equation}
with coefficients that depend on the chosen ratio or overlap-removal scheme.
Therefore the optional constant fit in $t_{\rm sep}$ should be interpreted only
as a large-separation consistency test unless a proper excited-state model is
fit.

As a first explicit excited-state model, the batch analyzer now also supports a
correlated one-exponential fit of the reduced fixed-$Q^2$ data,
\begin{equation}
  F_{\rm eff}(Q^2,t_{\rm sep})
  =
  F_\pi(Q^2)
  +
  A(Q^2)\,e^{-\Delta(Q^2)\,t_{\rm sep}}.
\end{equation}
This fit is performed after the $(Q^2,t_{\rm sep})$ channel averaging described
above, using the full jackknife covariance of the reduced data. It is a more
appropriate default reduced-$Q^2$ summary than a constant fit across all
$t_{\rm sep}$ values, but it is still only a first model. In particular, it
does not capture more complicated source/sink excited-state structures or
multiple gaps, so its fit quality must still be checked channel by channel.

To avoid obvious thermal contamination, the reduced-$Q^2$ fits now also apply
an automatic temporal-extent safety cut by default:
\begin{equation}
  t_{\rm sep} \le \frac{T}{2} - \Delta_T,
\end{equation}
with a configurable margin $\Delta_T$ (default $\Delta_T=1$). This removes the
most obvious wrap-around region from the reduced fits and records the excluded
large-$t_{\rm sep}$ values in the JSON output. This is a conservative filter,
not a replacement for a fully derived finite-$T$ fit ansatz.

\end{document}
